%=== CHAPTER TWO (2) ===
%=== Literature Review ===

\chapter{Literature Review}
\label{ch:lit_review}
\begin{spacing}{2.0}
%\setlength{\parskip}{0.2in}

\section{Introduction to Intelligent Transport Systems (ITS)}
The development of Intelligent Transport Systems (ITS) is necessary due to the expansion of cities worldwide. These systems investigate the traffic environment, particularly when congestion and accidents pose a risk to public safety. Traditionally, traffic is monitored using inductive loop detectors, sensors, and manual observation. These methods have significant drawbacks such as sensors that do not provide spatial context, fatigue-related mistakes and the inability to detect lateral interactions which lead to slower response times to incidents \cite{robles-serrano_automatic_2021}. To counteract these issues, deep learning and computer vision-based systems which utilise automated video analysis to provide a non-intrusive approach, have leveraged existing Closed-Circuit Television (CCTV) networks to gather detailed, high-dimensional vehicle paths, types, and interactions. The emergence of these systems has been facilitated by the advancement of Convolutional Neural Networks (CNNs) for object detection, alongside datasets such as the Car Accident Detection and Prediction (CADP) benchmark which acts as an essential resource for training and testing prediction models. Expanding on this, the primary capability of any vision-driven ITS starts with the identification and tracking of road users. This task requires a strict equilibrium between inference speed and accuracy, as models related to road and public safety must be both fast and accurate.

\section{Evolution of Object Detection Architectures}
Although initial two-stage detectors, such as Faster R-CNN, achieved high accuracy, they struggled with the latency demands of real-time monitoring, leading to the emergence of single-stage detectors like the You Only Look Once (YOLO) series \cite{zhang_research_2024, hussain_yolo-v1_2023}. Object detection for traffic surveillance requires a specific balance, which YOLO architectures provide by treating detection as a unified regression task \cite{zhang_research_2024}. This approach enables the direct mapping of image pixels to bounding box coordinates and class probabilities, avoiding the additional computational overhead associated with region proposal networks.

Due to the limitations of existing surveillance systems, transitioning from YOLOv5 to newer versions, such as YOLOv8 and YOLOv11, has become crucial \cite{li_deep_2025, hussain_yolo-v1_2023}. Among the beneficial new features is the incorporation of enhanced feature fusion techniques, such as Asymptotic Feature Pyramid Networks (AFPN), which enable models to gradually merge low-level texture details with high-level semantic data. These advancements can significantly enhance the detection of small objects \cite{li_deep_2025}. This is a vital improvement, as it facilitates the tracking of distant vehicles or pedestrians in wide-angle CCTV footage, where targets might occupy fewer than 30 pixels. Context-Guided (CG) modules and Dilated Reparameterization Blocks (DRP) are integrated into modern architectures; these successfully enlarge the network’s receptive field, allowing it to differentiate vehicles from background elements such as tree shadows or barriers \cite{hussain_yolo-v1_2023}. This enhancement in detection technology with strong temporal continuity to illustrate the ongoing trend of employing tracking-centric temporal analysis \cite{basheer_ahmed_real-time_2023}. This approach advances contextual insight into vehicle behaviour, elevating the system from a basic spatial localisation tool to a more comprehensive framework that identifies trajectories over successive frames.

\section{Advancements in Multi-Object Tracking}
Multi-object tracking has transitioned from foundational algorithms, such as Simple Online and Realtime Tracking (SORT), to more advanced systems capable of mitigating occlusion and camera instability, which are common issues encountered in outdoor traffic surveillance \cite{aharon_bot-sort_2022}.

Early versions employed Kalman filters \cite{kalman_new_1960} and the Hungarian algorithm \cite{kuhn_hungarian_1955}; however, these were not robust enough to handle identity switches during occlusions. This limitation was addressed by utilising deep appearance descriptors, such as DeepSORT; however, performance continued to suffer in dynamic outdoor settings where camera vibrations could resemble object movement. To tackle these problems whilst establishing a new benchmark for traffic monitoring performance, BoT-SORT (Bag of Tricks SORT) was developed \cite{aharon_bot-sort_2022}. Camera Motion Compensation (CMC) was incorporated into BoT-SORT to synchronise frame coordinates when camera vibrations occur \cite{aharon_bot-sort_2022}. A global homography matrix, which stabilises the reference frame, is subsequently generated using sparse optical flow. The Kalman filter is further enhanced by adjusting the state vector to incorporate the aspect ratio and height of the bounding box; this, in turn, allows the filter to effectively simulate the non-linear motion of manoeuvring vehicles. These enhancements and methodologies are crucial, since identity consistency is vital for anomaly detection systems that depend on continuous velocity profiles to identify abrupt halts or abnormal deceleration \cite{arriffin_vehicles_2023}. Beyond these developments, tracking and deriving vehicle speed from 2D video sequences presents a significant challenge, as it requires reconstructing 3D metric dynamics from the motion of projected 2D pixels.

\section{Vision-Based Vehicle Speed Estimation}
Determining vehicle speed is essential for understanding traffic flow and identifying anomalous driving behaviours. Nevertheless, it remains a considerable challenge within the domain of computer vision owing to the lack of depth information \cite{dong_vehicle_2019, hua_vehicle_2018, macko_efficient_2025}. Yang et al. \cite{dong_vehicle_2019} developed a deep learning model that utilises 3D CNN and non-local blocks to estimate vehicle speeds directly from video sequences. Their research demonstrated that integrating temporal context with spatial cues markedly improves accuracy compared to traditional frame-based estimation methods, which are frequently compromised by noise in bounding box placement. Similarly, in their vision-based framework, Hua et al. \cite{hua_vehicle_2018} utilised vehicle tracking and motion analysis to determine speed from traffic footage, confirming that observing the temporal displacement of a bounding box can serve as a highly reliable indicator of speed, provided the camera perspective is properly accounted for. In more challenging environments, such as night-time surveillance, a system was developed that tracks vehicles via their headlights, demonstrating that intelligent thresholding and preprocessing can maintain consistent detection even in low-light conditions \cite{hyungjun_vehicle_2019}. Through these approaches, recent research has increasingly focused on calibration-free methods that leverage the distribution of average vehicle sizes to convert pixel displacement into metric speed. This approach is particularly advantageous for large-scale deployment across urban areas, where the manual calibration of numerous cameras is impractical. A. Macko et al. have refined these techniques, achieving error margins as low as $\pm 0.58$ km/h on modern benchmarks by utilising 3D bounding box regression to determine scale from the vehicle’s known dimensions \cite{macko_efficient_2025}.

\section{Lane-Specific Analysis and Behaviour Monitoring}
This progress in speed estimation has significantly facilitated the transition from scene-wide traffic observation to a focus on individual lanes, shifting from aggregate vehicle counts to granular behavioural analysis. While overall speed estimation provides a collective overview of traffic movement, it can often obscure specific issues within lanes, such as a stalled vehicle in a turning lane while the remaining traffic flows smoothly. Lane-specific speed evaluation involves segmenting the road into individual lanes to maintain distinct speed profiles for each lane \cite{zhou_comprehensive_2022}. Consequently, irregularities can be identified more readily; for example, an automated system can flag a potential blockage in one lane while the others remain unobstructed. Furthermore, a localised approach enhances interpretability compared to previous models that solely provided aggregate traffic data \cite{arriffin_vehicles_2023}. Incorporating lane-association logic, which maps a vehicle’s centroid to a specific lane mask generated through semantic segmentation, is essential for this detailed analysis, as it enables the calculation of lane-specific metrics such as occupancy and flow rate. This granular data is integral to broader traffic anomaly detection systems, which are typically categorised into reconstruction-based, supervised classification, or hybrid rule-based methods.

\section{Traffic Anomaly Detection Frameworks}
Traffic anomaly detection primarily involves identifying events such as accidents and atypical movement patterns that deviate from standard traffic behaviour; this necessitates the capacity to distinguish between minor irregularities, such as vehicles momentarily stopping, and significant safety concerns.

A notable example of this methodology was proposed by Aboah et al. \cite{aboah_vision-based_2021}, who developed a hybrid system that combines YOLOv5 for spatial vehicle detection with decision tree logic for classifying anomalies. Their findings demonstrate that the integration of deep learning with rule-based reasoning significantly enhances both the interpretability and decision-making accuracy of the framework. This approach effectively addresses the interpretability challenges frequently associated with "black-box" models by explicitly decoupling visual perception from logical reasoning. The perceptual output from the neural network is subsequently analysed by an independent logical layer governed by predefined safety rules and physical constraints. This configuration facilitates the detection of specific anomalies, such as a vehicle remaining stationary for an extended period in an unsignalised area---defined as a zone lacking traffic control measures, such as a highway lane, where a halted vehicle clearly indicates an incident rather than routine traffic flow.

Furthermore, an alternative approach to spatial-temporal analysis was presented in the Intelligent Intersection framework, which employs a dual-stream Convolutional Neural Network (CNN) \cite{huang_intelligent_2019}. This architecture combines parallel appearance and motion streams to simultaneously aggregate both static visual features and dynamic kinematic data. This fusion allows the system to achieve greater sensitivity to complex vehicle interactions and imminent collision hazards. By continuously analysing raw pixel shifts alongside spatial tracking data, the CNN can detect early physical indicators of accidents, such as a vehicle skidding (atypical motion) prior to an impending structural impact (atypical appearance).

\section{Spatiotemporal Analysis via Dense Optical Flow}
Traditional object detection methods, whilst highly effective at isolating individual vehicle trajectories through bounding boxes, frequently encounter difficulties in interpreting the complex spatial dynamics and overall severity of multi-vehicle collisions \cite{robles-serrano_automatic_2021, aboah_vision-based_2021, huang_intelligent_2019}. When vehicles collide, their visual representations tend to intersect or fragment, rendering consistent temporal tracking highly complex. Recent scholarly work, such as that by Singh et al. \cite{singh_surveillance_2025}, emphasises the importance of capturing unstructured spatiotemporal dependencies to enable the accurate detection of road accidents. Although the authors proposed the utilisation of complex transformer-based architectures to model these temporal relationships, such approaches demand substantial computational overhead, which is frequently prohibitive for real-time deployment on edge devices within municipal surveillance contexts.

To quantify unstructured kinetic energy without the latency inherent to transformer models, numerous vision-based frameworks rely on dense optical flow techniques \cite{huang_intelligent_2019, goos_two-frame_2003}. Unlike sparse methods, such as the Lucas-Kanade algorithm, which track only a limited set of high-contrast corner points, the dense optical flow algorithm developed by Gunnar Farnebäck \cite{goos_two-frame_2003} calculates a unique movement vector for every individual pixel within the frame.

The mathematical foundation of the Farnebäck algorithm is based on quadratic polynomial expansion. The algorithm approximates the local image signal around each pixel in consecutive frames by modelling the intensity as a quadratic surface. For a specified 2D spatial coordinate vector, denoted as $\mathbf{x}$, the intensity within the initial frame is approximated as:

\begin{equation}
f_1(\mathbf{x}) \approx \mathbf{x}^T \mathbf{A}_1 \mathbf{x} + \mathbf{b}_1^T \mathbf{x} + c_1
\label{eq:farneback_poly1}
\end{equation}

Where:
\begin{itemize}
    \item $\mathbf{x}$ represents the 2D local spatial coordinate vector defined as $(x, y)^T$.
    \item $T$ denotes the standard matrix transpose operation.
    \item $\mathbf{A}_1$ is a symmetric matrix representing the second-order spatial derivatives, capturing the quadratic surface or curvature of the image intensity.
    \item $\mathbf{b}_1$ is a vector capturing the first-order derivatives, representing the linear gradient of the image intensity.
    \item $c_1$ is a scalar representing the constant term, corresponding to the absolute pixel intensity value at the centre of the local neighbourhood.
\end{itemize}

Assuming that the subsequent frame is subject to a global spatial translation $\mathbf{d}$ relative to the first frame, the polynomial for the second frame expands to:

\begin{equation}
f_2(\mathbf{x}) = f_1(\mathbf{x} - \mathbf{d}) = (\mathbf{x} - \mathbf{d})^T \mathbf{A}_1 (\mathbf{x} - \mathbf{d}) + \mathbf{b}_1^T (\mathbf{x} - \mathbf{d}) + c_1
\label{eq:farneback_poly2}
\end{equation}

By equating the coefficients of this transformed polynomial with the empirically observed coefficients in the second frame, the algorithm efficiently derives the precise displacement field, $\mathbf{d}$ \cite{goos_two-frame_2003}. 

In the context of traffic monitoring, extracting this raw kinetic signature enables tracking systems to map the exact location, direction, and magnitude of a collision as a spatiotemporal heatmap, irrespective of vehicle classifications or degraded environmental illumination.

\section{Deep Learning Classification using ResNet Architectures}
Following the extraction of kinetic features, the classification of anomaly severity necessitates the deployment of robust deep learning architectures \cite{robles-serrano_automatic_2021, basheer_ahmed_real-time_2023, aboah_vision-based_2021}. In the past, as CNNs got deeper to catch more complex details, they ran into the vanishing gradient problem. During the backward propagation of errors in training, the gradients diminished to negligible levels, resulting in the saturation and rapid deterioration of the network’s accuracy.

To overcome this architectural limitation, He et al. \cite{he_deep_2016} introduced the Deep Residual Network (ResNet) framework. The central innovation of ResNet lies in the residual block, which incorporates a "shortcut" or "skip" connection that bypasses one or more layers. Instead of forcing the stacked layers to learn a complex, unreferenced mapping $\mathcal{H}(\mathbf{x})$, the network is trained to learn a residual mapping $\mathcal{F}(\mathbf{x}) = \mathcal{H}(\mathbf{x}) - \mathbf{x}$ \cite{mahmood_new_2023}. The final output of the block thus becomes an additive identity mapping:

\begin{equation}
\mathbf{y} = \mathcal{F}(\mathbf{x}, \{W_i\}) + \mathbf{x}
\label{eq:resnet}
\end{equation}

This structure guarantees that gradients can propagate backwards through the network without obstruction, thereby facilitating substantially deeper architectures without compromising generalisation.

Within the field of Intelligent Transport Systems, ResNet architectures have demonstrated significant efficacy in interpreting traffic anomalies \cite{basheer_ahmed_real-time_2023, oluwadamilola_traffic_2026}. Recent work by Oluwadamilola and Emmanuel \cite{oluwadamilola_traffic_2026} highlighted the effectiveness of ResNet-18 in classifying the severity of traffic accidents. Their results suggest that ResNet-18 offers a preferable balance between classification accuracy and computational efficiency when compared to deeper, parameter-intensive architectures such as ResNet-50 or VGG-16. This makes it an ideal choice for rapid, real-time deployment at the edge.

Furthermore, to address the inherent class imbalances commonly encountered in traffic datasets, where normal driving cases significantly outnumber accidents \cite{shah_cadp_2018, lin_focal_2020}, classifiers such as ResNet are frequently combined with the Focal Loss function \cite{lin_focal_2020, oluwadamilola_traffic_2026}. Unlike standard Cross-Entropy loss, Focal Loss dynamically scales the penalty based on the model's confidence in its predictions:

\begin{equation}
FL(p_t) = -\alpha_t(1-p_t)^\gamma \log(p_t)
\label{eq:focal_loss}
\end{equation}

This formulation effectively down-weights the loss assigned to easy-to-classify examples (such as normal, steady traffic) and directs the network’s optimisation exclusively towards challenging, misclassified instances, such as differentiating between heavy braking and a low-speed collision.

\section{Model Ensembling for Enhanced Reliability}
To make these detections more reliable and mitigate the frequent false-positive rates observed in single-model systems, researchers are increasingly adopting model ensembling techniques. Ensembling combines predictions from different models to enhance overall accuracy and reduce false positives in surveillance footage \cite{natha_deep_2025}. Common ensemble methods, such as voting and stacking, are widely employed. In a voting ensemble, multiple classifiers (e.g., Random Forest, Support Vector Machines, CNNs) evaluate the same input data, with the final decision determined by a majority vote to reduce the bias inherent in individual models. Stacking offers a more advanced approach by training a meta-learner (typically a simple neural network or logistic regression) to synthesise the outputs of base models \cite{liu_hybrid_2025}. This technique often surpasses the performance of standalone models in assessing accident severity by learning which predictors are most significant in various situations. The I3D-ConvLSTM2D architecture exemplifies a more sophisticated multi-stream ensemble that combines RGB frames with optical flow information to identify both spatial features and motion irregularities over time \cite{adewopo_smart_2023}. This dual approach allows the model to fuse local object characteristics with the overall scene context, significantly improving Area Under the Curve (AUC) scores on standard datasets. Furthermore, recent studies have introduced "Expert Ensembles" (Xen), comprising specialised sub-models that monitor global flow, object interactions, and individual behaviour \cite{orlova_simplifying_2025}. These scores are subsequently combined using a Kalman filter to produce a stable, noise-resistant anomaly signal.

\section{Benchmark Datasets: The Role of CADP}
The efficacy of these advanced models is highly contingent upon the availability of robust benchmark datasets, with the Car Accident Detection and Prediction (CADP) dataset serving as a principal resource for the assessment of surveillance-based methodologies. CADP was created to address the scarcity of publicly accessible datasets containing precise spatio-temporal annotations pertinent to accident prediction \cite{shah_cadp_2018}. Comprising 1,416 video clips obtained from YouTube CCTV viewpoints, CADP presents standardised Time-to-Accident (TTA) annotations. The TTA metric is formally characterised as the temporal interval between an algorithm’s initial prediction of an anomaly and the precise moment of physical impact. These specific annotations, which on average precede a collision by 3.69 seconds, are fundamental to the development and evaluation of predictive early warning systems. This consistent temporal framework allows researchers to assess not only detection accuracy but also prediction latency, a factor of critical importance for collision avoidance. The third-person perspective embodied in CADP captures the comprehensive interactions among various traffic participants \cite{shah_cadp_2018}, rendering it markedly more advantageous than dashcam-centric datasets such as the Dashcam Accident Dataset (DAD) or AnAn Accident Detection (A3D) for monitoring urban infrastructure within smart cities. Whilst DAD and A3D furnish valuable ego-centric data relevant to autonomous vehicles, they often lack the broader spatial context provided by the stationary surveillance perspectives characteristic of CADP, which is essential for analysing flow dynamics across multiple lanes.

\section{Explainable AI and Edge Computing Integration}
Beyond these methodological and data-driven advancements, the integration of Explainable Artificial Intelligence (XAI) approaches addresses the "black-box" problem associated with deploying deep learning within safety-critical systems. Techniques such as Grad-CAM can generate heatmaps that highlight which parts of an image triggered an accident warning, such as a collision with an obstacle or an abrupt skid mark; thereby enabling operators to verify whether the model is attending to the relevant visual cues. In a hybrid framework, the rule-based layer is inherently transparent, as it explicitly indicates which kinematic rule has been violated (e.g., "Speed < 5 km/h in active lane"), which is essential for human operators requiring confirmation of automated alerts. XAI is vital for compliance with regulatory standards and for fostering public confidence in AI-driven monitoring systems, as it transforms a "black-box" alert system into a transparent decision-support tool \cite{pandey_explainable_2026}. To meet the demand for real-time detection, researchers are increasingly adopting edge computing architectures, whereby lightweight models such as Tiny YOLO are deployed directly onto camera devices. Localised processing reduces the transmission delays associated with cloud-based systems, thereby facilitating the sub-second response times necessary for accident prediction and emergency response \cite{somireddypalli_maneesh_kumar_ensemble_2025}. This architectural shift is often accompanied by model compression techniques, including quantisation and pruning, which enable complex ensemble models to operate efficiently on the resource-constrained hardware of edge devices.

\section{Multi-Modal Integration and Vision-Language Models}
Beyond relying solely on visual data, the future of traffic anomaly detection is increasingly dependent upon the integration of diverse data sources, such as LiDAR and V2X (Vehicle-to-Everything) communications. While camera-based systems are cost-effective, incorporating acoustic sensors or roadside LiDAR can significantly enhance accuracy, particularly in challenging weather conditions, such as heavy rain or snow, where optical sensors frequently encounter difficulties. Leveraging the geometric details of objects in video sequences can mitigate inherent visual biases in current datasets, thereby yielding models with improved generalisation capabilities \cite{mahmood_new_2023}. Furthermore, the emergence of Vision-Language Models (VLMs) offers a novel approach to traffic anomaly detection, substantially improving reliability within complex environments by enabling the semantic interpretation of traffic incidents. These models can describe an incident not merely as a set of coordinates, but as a coherent narrative (for example, "a white saloon failed to stop and collided with a pedestrian"), thereby streamlining communication with emergency services. Achieving this level of real-time comprehension requires advanced architectures that balance rapid reasoning with deep predictive capabilities, thereby expanding the overarching potential of automated traffic monitoring \cite{song_real-time_2025}.

\section{Conclusion}
In summary, prior research indicates that although individual components of traffic monitoring, such as YOLO for detection, BoT-SORT for tracking, and various methods for speed estimation have achieved notable progress, it is the integration of these technologies that holds the most promising prospects for the future. By synthesising the pattern recognition capabilities of deep learning with the deterministic logic of rule-based systems and the robustness of model ensembling, researchers can develop frameworks capable of comprehending the complex behavioural dynamics of traffic. The CADP dataset serves as a vital resource for these systems, addressing the scarcity of publicly available surveillance data despite its inherent limitations in detecting small objects. As the field progresses towards lane-level accuracy, improved interpretability, and real-time edge processing, the trend of developing hybrid systems to detect speed-based anomalies will remain pivotal in creating safer and more efficient Intelligent Transport Systems (ITS). This integrated approach not only facilitates proactive traffic management but also improves emergency response initiatives, ultimately contributing to the development of more resilient urban infrastructures that can enhance public safety on a broader scale.
%=== END OF CHAPTER TWO ===
\end{spacing}
\newpage
